\documentclass[11pt]{article}
\usepackage[margin=1in]{geometry}
\usepackage[T1]{fontenc}
\usepackage{lmodern,microtype}
\usepackage[colorlinks=true,allcolors=black]{hyperref}
\pagestyle{empty}
\begin{document}
\noindent Mathias Rodr\'iguez Castro\\
Facultad de Ingenier\'ia, Universidad de la Rep\'ublica\\
Montevideo, Uruguay\\
\href{mailto:mathiasr@fing.edu.uy}{mathiasr@fing.edu.uy}

\bigskip
\noindent To the Editors of \emph{Computational Optimization and Applications}

\bigskip
\noindent Dear Editors,

\medskip
Please consider ``Pre-Export Row Scaling in Generated Mixed-Integer Programs: Mechanism
Attribution and Residual-Budget Contracts'' for publication in \emph{Computational Optimization
and Applications}.

Positive row scaling preserves a mixed-integer program in exact arithmetic, but changes its
finite-precision behavior and the original-unit meaning of feasibility tolerances. Beyond
asking which scaling improves an optimization pipeline, this paper asks which mechanism
produces the apparent improvement and which accuracy guarantee survives the change of
representation.

Two analytical contributions organize the paper. First, geometric-mean centering eliminates
an extrema-based block statistic, making an apparently metadata-dependent coupling rule
identical to Euclidean normalization. The locally whitened triangular model admits the
closed-form optimal common coupling factor $(1+\rho^2)^{-1/2}$. Second, residual-budget
contracts characterize exactly the factors compatible with application-unit budgets and
coefficient limits, including an incompatibility certificate, constrained logarithmic minimax
centering, integer-rounding restrictions and conditions for exact binary64 export.
Kernel-matched controls distinguish the scaling mechanism from ownership metadata.

Computational validation includes a fixed-basis audit of 306 generated dispatch models,
1{,}080 matched matrix exports, 240 reconstructed operational runs with two matched
sensitivity reruns, and an exactly enumerated stress experiment with 4{,}752 optimizer calls
across three solvers. File-level conditioning gains largely shrink after solver-style
equilibration. A completed-run work reduction is reversed by timeout penalties and does not
reproduce with a second solver. Adding a residual budget to the same centering rule loses
no verified stress optima relative to that rule. These findings support numerical methodology
and error analysis; the operational evidence remains limited to one application.

The work connects the journal's areas of Integer Programming; Approximations and Error
Analysis; Software, Benchmarks, Numerical Experimentation and Comparisons; and Modelling Languages and Systems for
Optimization. Online Resource~1 provides proofs, protocols and reproduction instructions;
Online Resource~2 provides code, stored data, manifests and verification tools.

I am the sole author, received no external funding and declare no competing interests.
The manuscript documents AI assistance and my responsibility for its scientific content.
The work has not been published and is not under consideration elsewhere.

\bigskip
\noindent Sincerely,\\[6pt]
Mathias Rodr\'iguez Castro
\end{document}
